\uuid{RtEm}
\exo7id{7417}
\titre{exo7 7417}
\auteur{mourougane}
\organisation{exo7}
\datecreate{2021-08-10}
\isIndication{false}
\isCorrection{false}
\chapitre{Géométrie affine dans le plan et dans l'espace}
\sousChapitre{Géométrie affine dans le plan et dans l'espace}
\cpge{mpsi-mp2i.sous-espaces-affines, ptsi.espaces-vectoriels.b, mpsi-mp2i.applications-lineaires.e, pcsi.applications-lineaires.f, mpsi-mp2i.dimension-finie.c, pcsi.dimension-finie.c, ptsi.dimension-finie.c, bcpst1.geometrie.d, bcpst1.espace-vectoriel-kn.a, pt.complements-algebre-lineaire.d, bcpst2.espaces-vectoriels.a, bcpst2.applications-lineaires-matrices.b}
\module{Géométrie}
\niveau{L2}
\difficulte{}

\contenu{
\begin{enumerate}
    \item \question{Deux plans de l'espace affine $\Rr^3$ peuvent-ils avoir un unique point d'intersection ?}
    \item \question{Dans l'espace affine $\Rr^n$, quelle est la dimension 
d'un sous-espace affine défini par deux équations affines
indépendantes ?}
\end{enumerate}
}
